\chapter{Introduction}

\section{General Context}
In today's globalized economy, international trade and importation play a vital role in sustaining national markets. In Algeria, there is a continuously growing demand for imported goods driven by individuals and small-to-medium enterprises (SMEs) seeking products that are not locally available. The shift towards digital solutions in commerce has transformed how businesses operate, yet the importation sector in many developing regions still relies on traditional, manual, and often informal methods.

\section{Problem Statement}
Despite the high demand, the Algerian importation market remains highly fragmented, opaque, and complex for everyday users and professional importers alike. The current ecosystem is plagued by several major challenges:
\begin{itemize}
    \item \textbf{Lack of Centralization:} Importers and clients rely on scattered social media groups (Facebook, WhatsApp) and informal networks to conduct business, leading to a disorganized market.
    \item \textbf{High Risk of Fraud:} The absence of identity verification and secure payment gateways leaves clients vulnerable to scams.
    \item \textbf{Lack of Transparency:} Pricing, negotiation phases, and shipping fees are rarely transparent, causing disputes and a lack of trust.
    \item \textbf{Tracking Difficulties:} Clients have virtually no real-time visibility over their orders' shipment status once a transaction is initiated.
    \item \textbf{Barrier to Entry:} New entrepreneurs and individuals with small budgets face immense difficulties accessing micro-importation opportunities due to the lack of structured guidance.
\end{itemize}

\section{Proposed Solution: Sila DZ}
To address these critical shortcomings, we introduce \textbf{Sila DZ}, an intelligent, centralized B2B and B2C importation platform that directly connects clients with verified importers. 

Unlike conventional classifieds or social media groups, Sila DZ completely digitizes and structures the end-to-end importation lifecycle. The platform provides a highly secure environment where:
\begin{itemize}
    \item \textbf{Identity Verification:} All importers must undergo a rigorous identity check (using National ID cards and live camera verification) before they can operate, virtually eliminating fake accounts and fraudsters.
    \item \textbf{Integrated Negotiation:} A built-in, real-time negotiation engine allows clients and importers to discuss pricing, quantities, and terms in a structured chat, transforming informal haggling into a documented digital agreement.
    \item \textbf{Order Management \& Tracking:} Importers can publish detailed offers while clients can place custom orders. The system offers step-by-step real-time tracking, providing full visibility from the initial down payment to final delivery.
    \item \textbf{Business Model Innovation:} Sila DZ introduces a unique "Points System" allowing importers to purchase points to boost their offers' visibility or perform premium actions. This, combined with future subscription plans and transaction commissions, ensures a robust and scalable business model.
\end{itemize}

\section{Objectives}
The primary objective of this graduation project is to design, architect, and implement the Sila DZ platform from scratch using modern, scalable software engineering practices. Specific technical objectives include:
\begin{itemize}
    \item Architecting a robust, highly available Microservices backend.
    \item Implementing real-time communication systems for negotiations and notifications.
    \item Developing cross-platform interfaces (Web for dashboard/clients, Mobile for accessibility).
    \item Ensuring top-tier user experience (UI/UX) with built-in multi-language (RTL) support.
\end{itemize}

\section{Document Structure}
This thesis is structured into six chapters:
\begin{itemize}
    \item \textbf{Chapter 1: Introduction} outlines the context, problems, and the proposed solution.
    \item \textbf{Chapter 2: Background} provides fundamental definitions of the concepts and technologies driving modern distributed systems.
    \item \textbf{Chapter 3: Contribution} details the functional and non-functional requirements alongside the system's UML design.
    \item \textbf{Chapter 4: Implementation} discusses the software technologies, the microservices architecture, databases, and the development environments used.
    \item \textbf{Chapter 5: Result} showcases the final application, highlighting its features, strengths, and the value it adds to the users.
    \item \textbf{Chapter 6: Conclusion \& Perspectives} summarizes the work accomplished and proposes future enhancements.
\end{itemize}
